\documentclass[a4paper, 11pt]{article}

% -- Coments
\usepackage{verbatim}

% ----- Fonts -----
% -- Fuente --
% \usepackage{fontspec}
% \setmonofont{JetBrainsMono Nerd Font}  

% -- Color --
\usepackage{xcolor}
%\definecolor{azul}{RGB}{00,33,99}
\definecolor{azul}{RGB}{35,72,180}

% -- Page Margin --
\usepackage[margin=1in]{geometry}

% -- Espaciados --
\newcommand{\Pspace}{0.5cm}
\newcommand{\Aspace}{0.2cm}

% -- Columnas --
\usepackage{multicol}

% -- Imagenes --
\usepackage{graphicx}
\usepackage{float}

% -- Matemáticas --
\usepackage{amsmath, amssymb}

% -- Gráficas --
\usepackage{pgfplots}
\pgfplotsset{compat=1.18}

% -- Código --
\usepackage{listings}
\lstset{
    language=C++,                   % Lenguaje del código
    basicstyle=\ttfamily\small,     % Fuente del código
    keywordstyle=\color{blue},      % Color de palabras clave
    commentstyle=\color{gray},      % Color de comentarios
    stringstyle=\color{red},        % Color de cadenas
    numbers=left,                   % Números de línea a la izquierda
    numberstyle=\tiny\color{gray},
    breaklines=true,                % Permitir saltos de línea
    frame=single                    % Marco alrededor del código
}

\begin{comment}
-- Formato de pregunta simple
    % - Problema n
    \vspace{\Pspace}
    \item Problema \par
        % Respuesta:
        \vspace{\Aspace} \par
        { \color{azul}  }


 -- Formato de pregunta multimple
        % - Problema 
        \vspace{\Pspace}
        \item  \par
             % Respuestas:
            \vspace{\Aspace} \par
            a) 
            \\ { \color{azul} }

            \vspace{\Aspace} \par
            b) 
            \\ { \color{azul} }


 -- Formato de imagen
\begin{figure}[!ht]
    \centering
    \includegraphics[width=0.5\textwidth]{}
\end{figure}


 -- Formato de tabla
\begin{tabular}{c|c}
    \textbf{A} & \textbf{B} \\
    \hline
    Texto & Texto \\
    Texto & Texto
\end{tabular} \\
\end{comment}


\title
{
    Sistemas Operativos 2026-2 \\
    Calendarización en Sistemas Operativos de Tiempo Real
}

\begin{document}

\maketitle

\begin{center}
    \begin{tabular}{r|l}
        \textbf{Expediente} & \textbf{Nombre} \\ \hline
        219208106 & Bórquez Guerrero Angel Fernando \\
    \end{tabular}
\end{center}

\rule{\linewidth}{0.3mm}

\vspace{0.3cm}
\begin{enumerate}
    \item Dados los procesos descritos en la siguiente tabla: \par
    \begin{tabular}{|c|c|c|c|}
        \hline
        \textbf{Tarea} & \textbf{Periodo (p)} & \textbf{Periodo de CPU o CPU Burst (t)} & \textbf{Deadline (d)} \\
        \hline
        P1 & 50 & 35 & 50 \\
        \hline
        P2 & 80 & 25 & 80 \\
        \hline
    \end{tabular} \\
        % Respuestas:
        \vspace{\Aspace} \par
        a) ¿Cuál es la utilización de CPU?
        \par { \color{azul} 
        Para tareas periódicas, la utilización del CPU es $U=\sum_i \frac{t_i}{p_i}$:
        \[
            U = \frac{35}{50} + \frac{25}{80} = 0.70 + 0.3125 = 1.0125 \;\Rightarrow\; U = 101.25\,\%
        \]
        La utilización supera el 100\,\%, es decir, las tareas piden más tiempo de CPU del que el procesador puede ofrecer.
        }

        \vspace{\Aspace} \par
        b) Demuestre si se pueden calendarizar con el algortimo RMS.
        \par { \color{azul} 
        \textbf{Asignación de prioridades (RMS):} prioridad estática según el periodo; a menor periodo, mayor prioridad. Por lo tanto $P_1$ ($p=50$) tiene mayor prioridad que $P_2$ ($p=80$).

        \textbf{Prueba de utilización:} para $n=2$ tareas, RMS garantiza el cumplimiento de plazos si $U \le 2\left(2^{1/2}-1\right) \approx 0.8284$. Como $U = 1.0125 > 0.8284$ (y además $U>1$), no se cumple la condición.

        \textbf{Demostración con la línea de tiempo:}
        \begin{itemize}
            \item De $t=0$ a $t=35$ se ejecuta $P_1$ (termina antes de su plazo, $d=50$).
            \item De $t=35$ a $t=50$ se ejecuta $P_2$: alcanza a usar solo 15 de sus 25 unidades.
            \item En $t=50$ llega un nuevo trabajo de $P_1$, que tiene mayor prioridad y expulsa a $P_2$; $P_1$ se ejecuta de $t=50$ a $t=85$.
            \item En $t=80$ vence el plazo de $P_2$, y solo ha recibido 15 de sus 25 unidades (le faltan 10). $P_2$ reanuda en $t=85$ y termina en $t=95>80$.
        \end{itemize}
        \textbf{Conclusión:} las tareas \textbf{no} se pueden calendarizar con RMS, pues $P_2$ incumple su plazo en $t=80$.
        }
        
        \newpage
        \vspace{\Aspace} \par
        c) Demuestre si se pueden calendarizar con el algoritmo EDFS
        \par { \color{azul} 
        \textbf{Condición de EDF:} para tareas periódicas con plazo igual al periodo ($d=p$), EDF puede calendarizarlas si y solo si $U \le 1$. Aquí $U = 1.0125 > 1$, por lo que no se cumple.

        \textbf{Demostración con el hiperperiodo:} $\mathrm{mcm}(50,80)=400$. En 400 unidades de tiempo, $P_1$ se ejecuta $400/50=8$ veces y $P_2$ se ejecuta $400/80=5$ veces, de modo que el trabajo total es
        \[
            8(35) + 5(25) = 280 + 125 = 405 > 400 .
        \]
        El procesador solo dispone de 400 unidades, así que sobran 5 unidades de trabajo y, sin importar el algoritmo, algún plazo tiene que incumplirse.

        \textbf{Línea de tiempo:} EDF atiende siempre primero la tarea con el plazo más cercano. Cumple todos los plazos durante el primer tramo, pero la sobrecarga se va acumulando hasta que el trabajo de $P_1$ liberado en $t=350$ (plazo $d=400$) termina hasta $t=405$, es decir, 5 unidades tarde.

        \textbf{Conclusión:} las tareas \textbf{no} se pueden calendarizar con EDF. Cabe notar que EDF resiste más que RMS (el primer incumplimiento ocurre en $t=400$ y no en $t=80$), pero ambos fallan porque $U>1$.

        }



    \item Dados los procesos descritos en la siguiente tabla: \par
    \begin{tabular}{|c|c|c|c|}
        \hline
        \textbf{Tarea} & \textbf{Periodo (p)} & \textbf{Periodo de CPU o CPU Burst (t)} & \textbf{Deadline (d)} \\
        \hline
        P1 & 50 & 15 & 50 \\
        \hline
        P2 & 75 & 25 & 75 \\
        \hline
        P3 & 100 & 20 & 100 \\
        \hline
    \end{tabular} \\
        % Respuestas:
        \vspace{\Aspace} \par
        a) ¿Cuál es la utilización de CPU?
        \par { \color{azul} 
        \[
            U = \frac{15}{50} + \frac{25}{75} + \frac{20}{100} = 0.30 + 0.3333 + 0.20 = 0.8333 \;\Rightarrow\; U \approx 83.33\,\%
        \]
        El CPU está ocupado el 83.33\,\% del tiempo y queda libre el 16.67\,\% restante.
        }

        \vspace{\Aspace} \par
        b) Demuestre si se pueden calendarizar con el algortimo RMS.
        \par { \color{azul} 
        \textbf{Prioridades (RMS):} $P_1$ ($p=50$) $>$ $P_2$ ($p=75$) $>$ $P_3$ ($p=100$).

        \textbf{Prueba de utilización:} para $n=3$ tareas, el límite suficiente es $3\left(2^{1/3}-1\right) \approx 0.7798$. Como $U = 0.8333 > 0.7798$, la prueba no garantiza ni descarta la calendarización, por lo que se requiere un análisis exacto.

        \textbf{Análisis exacto (tiempo de respuesta en el peor caso):} $R_i = t_i + \sum_{j<i} \left\lceil \frac{R_i}{p_j} \right\rceil t_j$
        \begin{itemize}
            \item $R_1 = 15 \le 50$ \checkmark
            \item $R_2 = 25 + \left\lceil \frac{R_2}{50} \right\rceil 15$: con $R_2=40$ se cumple ($\lceil 40/50 \rceil = 1 \Rightarrow 25+15=40$), así que $R_2 = 40 \le 75$ \checkmark
            \item $R_3 = 20 + \left\lceil \frac{R_3}{50} \right\rceil 15 + \left\lceil \frac{R_3}{75} \right\rceil 25$: iterando desde $R_3=20$ se obtiene $60$, luego $75$, y se queda en $75$ (con $R_3=75$: $20 + 2(15) + 1(25) = 75$). Por lo tanto $R_3 = 75 \le 100$ \checkmark
        \end{itemize}
        \textbf{Línea de tiempo} (hiperperiodo $\mathrm{mcm}(50,75,100)=300$): el primer trabajo de $P_3$ es el más afectado, pues termina en $t=75$ (plazo 100); los siguientes terminan en $t=135$ (plazo 200) y $t=275$ (plazo 300). Todos los trabajos de $P_1$ y $P_2$ también cumplen sus plazos.

        \textbf{Conclusión:} las tareas \textbf{sí} se pueden calendarizar con RMS.

        }

        \newpage
        \vspace{\Aspace} \par
        c) Demuestre si se pueden calendarizar con el algoritmo EDFS
        \par { \color{azul} 
        \textbf{Condición de EDF:} con $d=p$, las tareas son calendarizables si y solo si $U \le 1$. Aquí $U = 0.8333 \le 1$, por lo que \textbf{sí se cumple}.

        \textbf{Verificación con el hiperperiodo (300):} $P_1$ se ejecuta 6 veces ($6 \times 15 = 90$), $P_2$ 4 veces ($4 \times 25 = 100$) y $P_3$ 3 veces ($3 \times 20 = 60$): el trabajo total es $250 \le 300$, y sobran 50 unidades en las que el CPU está libre.

        \textbf{Línea de tiempo:} EDF atiende siempre la tarea con el plazo más cercano. Todos los trabajos terminan antes de su plazo; por ejemplo, el primer trabajo de $P_3$ termina en $t=60$ (plazo 100), mejor que con RMS ($t=75$).

        \textbf{Conclusión:} las tareas \textbf{sí} se pueden calendarizar con EDF.
        }


\end{enumerate}
\end{document}
